\documentclass[12pt]{article}
\usepackage[utf8]{inputenc}
\usepackage[T1]{fontenc}
\usepackage{hyperref}
\title{Mixtec, Juxtlahuaca Language Resources for K-12 Curriculum Development}
\author{Compiled by Claude Sonnet 5 (claude-sonnet-5) for Andr\'as Kornai}
\date{}
\begin{document}
\maketitle
\section{Introduction}
Juxtlahuaca Mixtec (ISO 639-3 vmc, Glottocode juxt1235) is an Otomanguean language of Mexico, described as vigorous. The name is applied to more than one variety, so check which variety a resource covers.
\section{Primary Educational Resources}
No primers, textbooks or school programmes were found in search results.
\section{Linguistic and Reference Materials}
A 1979 Swadesh wordlist from Piedra Azul, Juxtlahuaca, with audio, is held at AILLA: \url{https://islandora-ailla.lib.utexas.edu/islandora/object/ailla%3A129133}. An online Juxtlahuaca Mixtec--English dictionary of mixed quality is at Glosbe: \url{https://glosbe.com/vmc/en}. Forum posts mention ILV dictionary and grammar materials, but these could not be verified and are not cited.
\section{Religious and Community Resources}
Global Recordings Network audio materials: \url{https://www.globalrecordings.net/en/language/vmc}.
\section{Digital and Audio Resources}
Wikitongues language page: \url{https://wikitongues.org/languages/vmc/}. Endangered Languages Catalogue entry for a related Southern Baja Mixtec listing: \url{https://endangeredlanguages.com/lang/10535?hl=es}.
\section{Curriculum Development Context}
The language is vigorous, which favours community-led, mother-tongue curriculum work. Published pedagogical materials were not found by this search and should be sought through Mexican education authorities and the community.
\end{document}